% Appendix: limitations in detail and licenses. Moved from Section~\ref{sec:limits}. Licenses were
% read on their official pages on 2026-10-30 (bib/STATUS.md).

\section{Limitations in detail}\label{sec:limits-limits}

\paragraph{Scope of the security claim.}
The title's \emph{leak-free} means that no returned chunk lay outside the readable set at the revision
the check read, under the trusted computing base of Section~\ref{sec:prob-criteria}. It does not mean
fresh: a revocation takes up to the $\delta$ of its path, which Section~\ref{sec:eng-revoke} measures,
and the release we evaluated exceeded its stated $\delta$ after an invalidation. It does not mean that
unreadable chunks leave no trace: latency, counts, recall and scores reveal them
(Section~\ref{sec:side}), and the detection fields of Section~\ref{sec:repair} (\emph{missed},
\emph{exhausted}, \emph{over\_limit}) are further caller-visible observables that we did not measure. A
store that binds content to the wrong identifier or metadata defeats any re-check, and we test that
assumption only as far as the conformance kit's metadata checks reach. The trusted computing base has
failed before: earlier versions of the layer's SpiceDB and Cedar adapters let the check itself admit
(Table~\ref{tab:conf-faults}), so the claim covers the code we evaluated, not its history. The
completeness flag is count-honest relative to the pushed-down filter: it detects under-retrieval but
not mis-retrieval, and a full answer with the wrong chunks still reports complete.

\paragraph{Synthetic permissions.}
No public corpus pairs documents with the access-control lists of a real organization, so every
permission layout here is generated: tenants, groups and per-document ACLs drawn from templates and
seeded generators, over public passages (FiQA, MS MARCO) and templated company documents. Real
organizations may have more skewed layouts, with many small groups, deep group hierarchies and
frequent exceptions. We place every (question, subject) pair on measured correlation axes instead of
trusting a layout's label (Section~\ref{sec:prob-metrics}), which makes results comparable across
layouts of similar measured difficulty, although no single axis described every store
(Section~\ref{sec:qual-frontiers}). The MS MARCO experiments vary layouts across seeds; the
end-to-end runs keep one layout and vary only documents and questions, so their people recur with the
same permissions. That one layout is harsh in a way we state and did not vary: each company document
has 19 near-identical copies in other tenants that only the tenant rule separates, a hard case for
graph indexes, so its recall numbers, Qdrant's especially, should not be read as an estimate for an
organization without such copies.

\paragraph{One instrument.}
Every measurement passes through one enforcement layer, built for this project, so a fault in its
filter translators could be mistaken for a store fault. The conformance kit compares every answer
with an in-memory reference that shares no code with the translators, and the retrieval audit uses
each store's own exact search as a second baseline. Some numbers are properties of the layer's schema
as much as of a store: pgvector's filters are \texttt{jsonb} containment tests over a GIN index, with an
array test above 100 values that the index cannot serve, and Chroma's list values are boolean keys, so
its disjunctions hold one clause per group. Results describe each store under this layer's schema, not
at its best possible schema. Of the third-party behaviors our experiments found, four come with a
reproducer that uses only the vendor's own client: the truncation in Redis and in MongoDB, Chroma's
crash, and the CEL semantics that the layer's Cerbos translator meets. Each showed the reported
behavior when run on GitHub-hosted runners on 30 October 2026 against the version we evaluated and the
vendor's latest release (Redis 8.10.2 and 8.10.3, MongoDB Atlas Local 8.0.30 and 8.3.13, Cerbos 0.56.0
and 0.57.0, and Chroma 1.5.9, still its latest). Running them corrected one figure: MongoDB's
vector-search filter fields keep 8,191 characters, not the 8,181 we had taken from its documentation
of another field type. SpiceDB's staleness, OPA's residual negation and pgvector's plan choice were
measured through the layer and have no reproducer independent of it; the other rows of
Table~\ref{tab:conf-faults} rest on the vendors' documentation, which the table cites.

\paragraph{Versions, settings and hardware.}
No single version of the library was evaluated throughout. The end-to-end runs used release 2.1.0 from
PyPI; other experiments ran from their own branches, several carrying fixes that are not yet in a
release, such as the engine rerun and the repair experiment (Appendix~\ref{app:map}), and we did not rerun the whole evaluation
on one release. The library commit of each experiment is recorded with its outputs. Results hold for the
versions and settings we state (for example pgvector 0.8.7 and Qdrant 1.19.2) and may not survive the
next release; vector stores change quickly, and several behaviors we measure are documented defaults
that vendors may revise. PostgreSQL was 16.15 in the end-to-end runs and version 16 in the diagnosis
grid, and 17.11 in the campaign, the engine study and the frontier, repair and filter-size experiments.
Embeddings were 768-dimensional (bge-base-en-v1.5) on the enterprise corpus and 384-dimensional on MS
MARCO (bge-small-en-v1.5) and in the filter-size and side-channel experiments (all-MiniLM-L6-v2);
dimension matters, for example, to the point count at which Qdrant plans an exact scan. We did not
record Lucene's segment counts for Elasticsearch and OpenSearch. Experiments ran on shared cloud
hardware (Kaggle sessions with two T4 GPUs and four CPUs, and GitHub-hosted runners), so latencies are
comparable within an experiment, not across experiments or with vendor benchmarks.

\paragraph{Baselines and sweeps we did not run.}
We compare the enforcement paths of one layer, not the layer with a store's native enforcement: we
did not measure PostgreSQL row-level security or the document-level security of Elasticsearch and
OpenSearch, whose recall and latency may differ. We did not sweep the share of near-duplicate
chunks, so we cannot say how fast Qdrant's loss falls as copies across tenants become fewer. Outside the
100\% selectivity cells on MS MARCO we have no unfiltered baseline per store, so part of a store's loss
under a filter may be its ordinary approximation error, as with LanceDB's quantized index.

\paragraph{Statistical power.}
The end-to-end scenario ran four times with fresh documents and questions, and four more times for
the OpenFGA configuration, but always with one permission layout; its people recur in every run, so
its bound per person is a single run's. Its bounds are over the question generator's distribution and
say nothing about an adaptive adversary. The diagnosis grid ran three index builds per layout, the
frontier and repair experiments three builds per store and corpus over ten MS MARCO layouts and the
enterprise layout, and the engine, filter-size and side-channel experiments five seeds each; intervals
come from bootstrap resampling over people, subjects or probes, not over builds, so a store whose
recall varies between builds (Qdrant's by up to 0.13) is described by the mean of three. Zero-leak
results are reported as upper bounds per question, per subject and per fault episode, not per hit. We
apply no correction for multiple comparisons across cells, and in the side-channel experiment two
builds of identical content can differ by as much as a weak channel (Section~\ref{sec:side-results}).
Several pre-registered verdicts are inconclusive rather than confirmed or falsified, and we report
them as such.

\subsection{Licenses of data and models}\label{app:licenses}

We read each license on its official page on 30 October 2026. MS MARCO's terms state that the datasets
are intended for non-commercial research purposes only. The FiQA site makes its Task 2 data, which
BEIR redistributes, available only for non-commercial use; BEIR's dataset cards for FiQA and SciDocs
carry a CC BY-SA 4.0 tag, and BEIR's documentation leaves it to users to determine whether they may use
a dataset. The SciDocs repository's license file is CC BY 4.0. The embedding models bge-base-en-v1.5 and
bge-small-en-v1.5 are under the MIT license and all-MiniLM-L6-v2 under Apache 2.0. Qwen2.5-3B-Instruct
is under the Qwen Research License Agreement, which permits research and evaluation and requires a
separate license for commercial use. All our use was research. The harnesses download corpora and
models from their sources; released outputs that quote passages or model answers should be treated
under the same terms.
